\documentclass[11pt]{article}
\usepackage[T1]{fontenc}
\usepackage{lmodern}
\usepackage[margin=1in]{geometry}
\usepackage{amsmath,amssymb,amsthm,mathtools}
\usepackage{microtype}
\usepackage{enumitem}
\usepackage{xcolor}
\usepackage{longtable,booktabs,array}
\usepackage[colorlinks=true,linkcolor=blue!45!black,citecolor=blue!45!black,urlcolor=blue!45!black]{hyperref}
\usepackage{url}
\urlstyle{same}
\hypersetup{pdftitle={Factor-count residues and Poisson-Dirichlet limits for prime predecessors},pdfauthor={Myeongjin Kang}}
\setlength{\emergencystretch}{2em}
\setcounter{tocdepth}{1}
\numberwithin{equation}{section}
\newtheorem{theorem}{Theorem}[section]
\newtheorem{lemma}[theorem]{Lemma}
\newtheorem{proposition}[theorem]{Proposition}
\newtheorem{corollary}[theorem]{Corollary}
\theoremstyle{definition}
\newtheorem{definition}[theorem]{Definition}
\newtheorem{assumption}[theorem]{Assumption}
\theoremstyle{remark}
\newtheorem{remark}[theorem]{Remark}
\newcommand{\N}{\mathbb N}
\newcommand{\Z}{\mathbb Z}
\newcommand{\R}{\mathbb R}
\newcommand{\C}{\mathbb C}
\newcommand{\E}{\mathbb E}
\newcommand{\Prob}{\mathbb P}
\newcommand{\ind}{\mathbf 1}
\newcommand{\eps}{\varepsilon}
\newcommand{\dd}{\,\mathrm d}
\newcommand{\W}{\mathcal W}
\newcommand{\Bx}{\mathcal B_x}
\newcommand{\PD}{\operatorname{PD}}
\newcommand{\GEM}{\operatorname{GEM}}
\newcommand{\e}{\mathrm e}
\DeclareMathOperator{\lcm}{lcm}
\DeclareMathOperator{\Rea}{Re}
\DeclareMathOperator{\Ima}{Im}
\DeclareMathOperator{\Var}{Var}
\title{Factor-count residues and Poisson--Dirichlet limits\newline for prime predecessors}
\author{Myeongjin Kang\\\small Independent researcher}
\date{October 9, 2026\\[3pt]\small Working draft for external review}
\begin{document}
\maketitle
\begin{abstract}
Let $p$ range over the primes in a fixed reduced arithmetic progression.
We give a proof, conditional on specified analytic results in the OpenAI
mathematics repository, that $\Omega(p-1)$ modulo any fixed positive integer
is asymptotically uniform and independent of the ranked logarithmic sizes
of the prime factors of $p-1$. The limiting factor-size distribution is
$\PD(1)$. For parity, the joint limit also includes the squarefree indicator,
with its explicit progression-dependent density. This implies the
Moree--Hommersom parity conjecture under the same assumptions.
The proof treats the Liouville sign by an uncentered marked correlation.
For nonreal roots of unity, it constructs a common principal term from the
full small-prime part of each integer. This term respects shared prime
labels and can be summed through the allocation identities before its
absolute value is estimated. The assumptions refer to an immutable source
revision and are listed by statement in the paper. The draft was developed
with AI assistance and has not undergone external human peer review or
formal proof verification.
\end{abstract}
\tableofcontents
\clearpage
\section{Statements and assumptions}\label{sec:statements}

For $n\geq1$, let $\Omega(n)$ count prime factors with multiplicity.
For each prime $p\geq3$, list the prime factors of $p-1$ in decreasing
order, with multiplicity, as $q_1(p),q_2(p),\ldots$, and put
\[
 V_i(p)=\frac{\log q_i(p)}{\log(p-1)}.
\]
After the last prime factor, set $V_i(p)=0$. The sequence $V(p)$ is a
decreasing partition of one. Let $U_1,U_2,\ldots$ be independent uniform
variables on $(0,1)$, and define
\[
 T_i=U_i\prod_{j<i}(1-U_j).
\]
The sequence $(T_i)$ has the $\GEM(1)$ distribution. Its decreasing
rearrangement $L=(L_1,L_2,\ldots)$ has the Poisson--Dirichlet distribution
$\PD(1)$. We write $\pi(X;M,b)$ for the number of primes at most $X$
congruent to $b$ modulo $M$.

The distribution of $V(p)$ alone does not determine the parity of
$\Omega(p-1)$, or its residue modulo a larger integer. Small prime factors
can change that residue without producing a detectable change in finitely
many normalized factor sizes. We therefore prove cancellation of each
nontrivial root-of-unity weight against marked factor statistics before
passing to a probability limit.

\begin{assumption}[Analytic source inputs]\label{ass:inputs}
Assume the specified results and proof estimates in the following
manuscripts, at revision
\texttt{fd4aeeb2ee4fc729c18d98444fed42fd0529eeeb} of the OpenAI
mathematics repository:
\begin{enumerate}[label=\textup{(\roman*)}]
\item the graph transference and residual ideal-family inputs, comparison
estimates, long-prime estimates, and positive extraction estimates in
\emph{The Poisson--Dirichlet law for prime predecessors} \cite{D};
\item the general physical pairing and ideal-kernel results in
\emph{Weighted dilation graphs, smooth shifted primes, and totient fibers}
\cite{S}, as restated in \cite[Section 6]{D};
\item the scalar coefficient criterion in
\emph{Primitive roots for every admissible integer base}
\cite[Section 8]{P}, including the scalar Cauchy-square estimate in its proof;
\item the zero-free half-plane $\Rea s>7/8$ for Dirichlet $L$-functions
in \emph{The quasi-Riemann hypothesis} \cite{Z}, with the principal pole
at $s=1$ retained.
\end{enumerate}
The precise source labels, hypotheses needed here, and uses are recorded
in Section~\ref{sec:sources}. These source results are premises of this paper.
\end{assumption}

The fixed-progression Poisson--Dirichlet law is proved below as
Proposition~\ref{prop:unsigned}; it is not inferred by conditioning the
unrestricted source law. Its progression argument is an explicit
dependency of both the residue theorem and the squarefree refinement.

\begin{theorem}[Factor-count residues]\label{thm:main}
Assume Assumption~\ref{ass:inputs}. Fix integers $M,m,k\geq1$, a reduced
class $b\pmod M$, a residue $a\pmod m$, and a continuous function
$H:[0,1]^k\to\C$. Then
\begin{equation}\label{eq:main}
 \frac{1}{\pi(X;M,b)}
 \sum_{\substack{3\leq p\leq X\\p\equiv b\ (M)}}
 \ind_{\Omega(p-1)\equiv a\ (m)}H(V_1(p),\ldots,V_k(p))
 \longrightarrow \frac1m\E H(L_1,\ldots,L_k).
\end{equation}
The same limit holds for primes in $(x,2x]$, with denominator
$\pi(2x;M,b)-\pi(x;M,b)$.
\end{theorem}

Every displayed parameter is fixed before the real endpoint tends to
infinity. The theorem gives no rate and no uniformity in growing moduli.
Taking $H=1$ gives uniform residues; dividing by their limiting
probability $1/m$ gives the conditional factor-size law.

Write $\lambda(n)=(-1)^{\Omega(n)}$ and let $\mu$ be the M\"obius
function. For a prime $q$, set $e_q=v_q(M)$ and define
\begin{equation}\label{eq:rho}
 \rho_q(M,b)=
 \begin{cases}
  1/[q(q-1)],&e_q=0,\\
  1/q,&e_q=1,\ b\equiv1\pmod q,\\
  0,&e_q=1,\ b\not\equiv1\pmod q,\\
  1,&e_q\geq2,\ b\equiv1\pmod{q^2},\\
  0,&e_q\geq2,\ b\not\equiv1\pmod{q^2}.
 \end{cases}
 \qquad C_{M,b}=\prod_q(1-\rho_q(M,b)).
\end{equation}
The product includes $q=2$.

\begin{theorem}[Parity and squarefreeness]\label{thm:squarefree}
Under Assumption~\ref{ass:inputs}, if $p\geq3$ is uniform among the class
primes in either sampling convention of Theorem~\ref{thm:main}, then
\[
 (V_1(p),\ldots,V_k(p),\lambda(p-1),\mu^2(p-1))
 \Longrightarrow\ (L_1,\ldots,L_k,S,B),
\]
where $S$ is a fair sign, $B$ is Bernoulli with parameter $C_{M,b}$,
and $L,S,B$ are mutually independent. In particular, for $s\in\{-1,1\}$,
\begin{align}
 \E[\mu(p-1)H(V_1(p),\ldots,V_k(p))]&\longrightarrow0,\label{eq:mu}\\
 \E[\ind_{\mu(p-1)=s}H(V_1(p),\ldots,V_k(p))]
   &\longrightarrow\frac{C_{M,b}}2\E H(L_1,\ldots,L_k).\notag
\end{align}
\end{theorem}

The notation $\E$ on the left denotes the finite prime average.
If a square is forced by the class, $C_{M,b}=0$ and the Bernoulli
variable is zero. Conditioning on squarefreeness or on a M\"obius sign
requires $C_{M,b}>0$. For $M=1$, each M\"obius sign has density
$\tfrac12\prod_q(1-1/[q(q-1)])$ relative to the primes.

For a finite set $\mathcal S$ of primes, define
\[
 \mu_{\mathcal S}(n)=\mu\left(n\Big/\prod_{q\in\mathcal S}q^{v_q(n)}\right).
\]
Thus square factors in $\mathcal S$ do not force this function to vanish.

\begin{corollary}[Prescribed valuations]\label{cor:moree}
Under Assumption~\ref{ass:inputs}, for every finite $\mathcal S$ and
fixed nonnegative integers $e_q$, $q\in\mathcal S$,
\[
 \sum_{\substack{p\leq X\\v_q(p-1)=e_q\ (q\in\mathcal S)}}
      \mu_{\mathcal S}(p-1)=o(\pi(X)).
\]
This is the parity assertion of Moree and Hommersom
\cite[Conjecture 1]{MH}, under the stated source assumptions.
\end{corollary}

\subsection{Related work and the proof's scope}

The ordinary-integer distribution results of Pintz \cite{Pintz} and
Granville and Shao \cite{GS} supply progression estimates used below.
They do not restrict the averaging variable to prime predecessors.
Lichtman \cite{Lichtman} proves cancellation averaged over a growing
range of shifts. Carella \cite[Theorems 1.1--1.2]{Carella} states
shifted-prime cancellation subject to a separate distribution hypothesis.
These statements must be distinguished from the fixed-shift joint law
in \eqref{eq:main}. We make no claim to an exhaustive priority search.

The factor-count residue is recovered by finite Fourier inversion.
The analytic problem is to prove cancellation for
$f_z(n)=z^{\Omega(n)}$, for each nontrivial root of unity $z$.
For $z=-1$, the ordinary mean has arbitrary logarithmic cancellation.
For nonreal $z$, the Selberg--Delange expansion has a nonzero principal
term. Sections~\ref{sec:model} and~\ref{sec:extraction} construct and
propagate that term. Section~\ref{sec:probability} derives the joint
laws from the resulting interior factor estimates.

\section{Marked sums and analytic inputs}\label{sec:framework}

Put $L=\log x$ and $\e(t)=\exp(2\pi it)$. A constant denoted by
$C$ may change between estimates. Every exponent of $L$ is fixed before
$x$ tends to infinity. The notation $O_A(L^{-A})$ means that each fixed
accuracy $A>0$ is available, after the stated choices of parameters.

\subsection{Long slots, inert bands, and active groups}

Fix $d\geq1$ and $\eps>0$. Let $I_1,\ldots,I_d$ be intervals of
primes whose lower endpoints are at least $x^\eps$ and whose upper
endpoints have product at most $x^{1-\eps}$. Tuples are ordered and
their numerical prime entries are distinct. Write
\begin{equation}\label{eq:tuple}
 T_x(n)=\sum_{\boldsymbol\ell}\ind_{\ell\mid n},\qquad
 \ell=\ell_1\cdots\ell_d,\qquad C_x=\sum_{\boldsymbol\ell}\frac1\ell.
\end{equation}
On each fixed range $n\leq x^C$, $T_x(n)=O_{C,d,\eps}(1)$.
Also $C_x=O_{d,\eps}(1)$ by the harmonic prime estimate.

An inert array consists of $K$ disjoint full prime intervals
$[\exp(L^{a_i}),\exp(2L^{a_i}))$, with fixed $.1<a_i<.2$.
For damping $0<q<1$, its weight is
\begin{equation}\label{eq:inert}
 \W(n)=\prod_{i=1}^K q^{\omega_i(n)-1}\frac{\omega_i(n)}{V_i},
 \quad \omega_i(n)=\sum_{p\in\mathcal P_i}\ind_{p\mid n},
 \quad V_i=\sum_{p\in\mathcal P_i}\frac1p.
\end{equation}
A factor is zero when $\omega_i(n)=0$. For fixed $K,q$, the weight
is bounded. Its independent divisor-model mean $m_{\W}$ is obtained
by replacing the divisibility indicators by independent Bernoulli
variables with means $1/p$. The source gives
$0<c_K\leq m_{\W}\leq C_K$ and $\E_{\rm div}\W^2\leq C_K$,
uniformly over the fixed exponents.

Independently, choose disjoint active groups $\mathcal P_g$, in two
families $\mathcal G_s,\mathcal G_b$. Each family has between fixed
positive multiples of $\log L$ groups, and
\begin{align*}
 \mathcal P_g&\subset[\exp(L^{.27}),\exp(L^{.36})],&&g\in\mathcal G_s,\\
 \mathcal P_g&\subset[\exp(L^{.39}),\exp(L^{.46})],&&g\in\mathcal G_b.
\end{align*}
Their harmonic masses lie in a fixed compact subinterval of $(0,\infty)$.
Every big group is a full prime interval. At least one small group is
also a full interval, as needed for fixed-modulus mixing. The extraction
uses full intervals for all candidate groups.

Let $n_*$ denote $n$ after removal of all powers of active primes.
For $\mathbf t=(t_g)$ and a bounded function $b$ of ordered distinct
represented divisor lists, set
\begin{equation}\label{eq:active}
 W_{\mathbf t}^{b}(n)=
 q_0^{\omega_P(n)-\sum_g t_g}\prod_g V_g^{-t_g}
 \sum_{\text{represented lists at }n} b(\mathbf p),
 \qquad 0<q_0<1.
\end{equation}
Here $\omega_P(n)=\sum_g\omega_g(n)$ counts distinct active prime
divisors, and $V_g=\sum_{p\in\mathcal P_g}1/p$.
Write $W_1$ when $t_g=1$ for all $g$ and $b=1$.
For fixed bounded list lengths and $|b|\leq1$, these weights are at
most a fixed power of $L$. The second endpoint will have the form
$G(n)=G(n_*)$ with $|G(n)|\leq L^C\tau(n_*)^C$.

\subsection{The source estimates used at the two endpoints}

We use the general physical pairing in \cite[Section 6]{D}, rather
than extending its final centered correlation theorem by substitution.
The following description fixes its measures and the parts of its
conclusion that enter the proof.

For a sufficiently large fixed integer $J$, a physical state at $n$
has $J+1$ ordered distinct prime divisors from each active group. Its
list mass is $\prod_g V_g^{-(J+1)}$, with counting measure in $n$.
In a pattern $\nu$, the source and target share $J-|C_g|$ labels in
group $g$. Small groups have $C_g=\varnothing$; big groups have
$|C_g|\leq m_0$ for a fixed $m_0$. Let $D_R$ be the shared product
and $D_C$ the product of the auxiliary free labels. Free labels have
independent laws $1/(V_gp)$. The physical shift is $D_RD_C$.
The pattern kernels depend only on free big labels and unshared big
endpoint lists. The family has total kernel mass $L^{O(1)}$, with
exponent independent of $J$.

The residual ideal-family theorem supplies a raw pattern $D_C=1$
and signed comparison patterns. Its ideal operator has arbitrarily
small logarithmic norm. The general transfer theorem then bounds the
whole physical pairing for a first vector independent of its lists at
fixed $n$, with the source supremum bound, and endpoint norms
$\leq X^{1/2}L^{C_3}$ on a position dyad $[X,2X)$.
The exponent $C_3$ must be independent of $J$.
Both endpoints retain the extra half damping
$q_0^{(\omega_P(n)-(J+1)|P|)/2}$, which is at most one on represented
states. If the first core is bounded and the second has the divisor
bound above, the required norm estimate follows from
\begin{equation}\label{eq:listnorm}
 \prod_g V_g^{-(J+1)}\sum_{\text{lists}}\frac1{\prod p}\leq1
\end{equation}
and a fixed divisor moment. Conjugate the first core in the physical
vector to cancel the conjugation in the Hilbert-space pairing.

After inserting $\Psi(n/(xD_RD_C))/n$, restoring its full Mellin
integral, and summing all position dyads, put
\begin{equation}\label{eq:strip}
 n=D_Rw,\quad n'=D_R(w+D_C),\qquad
 \frac{\Psi(n/(xD_RD_C))}{n}
   =\frac1{D_Rw}\Psi\left(\frac{w}{xD_C}\right).
\end{equation}
The remaining cutoff has no shared-label dependence. The factor
$D_R^{-1}$ makes the shared-label measures harmonic probabilities.
The source's collision and Mellin-tail errors save every fixed power
of $L$, also after bounded multipliers at the endpoints.

Every comparison has two nonempty free big products, at scales
$U_1,U_2$. Write
\begin{equation}\label{eq:scales}
 H'=U_1U_2,\quad Y=xH',\quad H=H'/L^{C_6},\quad
 T_0=LY/H=xL^{C_6+1},\qquad
 U_i\geq\exp(L^{.39}-O(1)).
\end{equation}
Expanding the source kernels reduces a comparison to bounded separate
tests on the two free tuples, with a total fixed logarithmic cost.
Its additive multiplier is small off rational arcs of denominator
$\leq L^{C_6}$ and width $H^{-1}$. The source's local Mellin estimate
reduces the first-endpoint energy on these arcs to
\begin{equation}\label{eq:energy}
 \int_{|t|\leq T_0}
 \left|Y^{-1}\sum_n \mathcal U(n)\e(hn/k)n^{it}\right|^2\dd t
 \ll_A L^{-A},\qquad k\leq L^{C_6}.
\end{equation}
Here $\mathcal U$ has a smooth cutoff at $Y$, a logarithmic-sized
twist, the proposed first core, and a factor
$\W W_{\mathbf t}^{b}$ whose test $b$ uses only big marks.
The second endpoint needs only $\sum|\mathcal V(n)|^2\ll YL^C$.
The estimates and kernel expansions in this paragraph are the input
package from \cite[Section 6]{D}; we prove \eqref{eq:energy} for the
new cores below.

The truncated divisor model for $\W W_{\mathbf t}^{b}$ has the form
\begin{equation}\label{eq:divisor}
 \sum_{a\mid n}A_a+E(n),\quad
 a\leq\exp(O(L^{.47})),\quad
 \sum_a\frac{|A_a|}{a}\leq L^C,\quad
 \sum_{n\asymp R}|E(n)|^2\ll_N RL^{-N}.
\end{equation}
It permits a zero list length when a small mark has been removed.
The constants are uniform in bounded big-mark tests.

We also use two explicit inputs from \cite[Appendix 9]{D}.
For fixed $0<\tau<\eta<1$ and $A,C>0$, there is $B_0$ such that
\begin{equation}\label{eq:longprime}
 \left|\sum_{p\in I}\chi(p)p^{-1+it}\right|\ll L^{-A}
\end{equation}
for $I\subset[N,2N]$, $x^\tau/2\leq N\leq x^\eta$,
$\operatorname{mod}(\chi)\leq L^C$, and $L^{B_0}\leq|t|\leq x^2$.
For some absolute $C_*$, the second estimate is
\begin{equation}\label{eq:logphase}
 \left|\sum_{\substack{n\in I\\n\equiv a\ (k)}}n^{it}\right|
 \ll\frac{N'}k\exp\left(-\frac{L}{(\log L)^{C_*}}\right),
\end{equation}
where $I\subset[N',2N']$, $k\leq L^C$,
$\exp(L/(\log L)^2)\leq N'\leq2x^5$, and
$\exp(L/(2(\log L)^2))\leq|t|\leq4x^3$.
The residue $a$ need not be a unit. The subinterval uniformity in both
estimates allows partial summation against the smooth amplitudes below.

\subsection{The scalar determinant and positive extraction inputs}

The coefficient criterion in \cite[Section 8]{P} applies to rough
dyadic coefficients $\alpha_m,\beta_n$, bounded by fixed powers of $L$,
on positive-power scales $H_mH_n\asymp x$. It requires
\begin{equation}\label{eq:coefficient}
 H_m^{-1}\left|\sum_m\alpha_m\chi(m)m^{it}\right|\ll_A L^{-A}
\end{equation}
for every fixed logarithmic modulus range, every fixed accuracy, and
$|t|\leq2H_mH_nL^B$. The scalar estimate in its proof bounds the
expanded Cauchy square after $mn-1=ah$, with $a$ a product of $K$
selected inert primes. An invariant bounded factor of $h$ lies outside
that square. The source prime-minus-rough coefficient meets
\eqref{eq:coefficient}. Multiplication by a character of a fixed
modulus $M$ replaces the test modulus by its least common multiple
with $M$, which is still logarithmic.

The extraction in \cite[Section 7]{D} uses a nonnegative block sieve,
prime-to-rough telescoping, and a partition into independent conditional
coin outcomes. The block sieve has relative error $O(e^{-h})$ at even
depth $h$ and uses summed progression remainders only up to
$z^{4h+2}$ when the sieve limit is $z$. For the independent divisor
model, its local densities are $1/r$ at allowed primes.
The rough proxy is
\[
 a(m)=\frac{\ind_{m>x^\gamma}\ind_{P^-(m)>W}}{V(W)\log m},
 \qquad W=\exp(L^{.24}),\quad V(W)=\prod_{p\leq W}(1-1/p).
\]
The source permits removal of all candidate prime parts from its
thresholds and logarithms at absolute cost $o(x/L)$. It also gives a
total positive proxy mass $O_{j,\gamma}(x/L)$ for each fixed number
$j$ of factors. Section~\ref{sec:extraction} states the allocation
identity and carries out the needed new recombination.
These imported estimates do not assert cancellation of $f_z(p-1)$.

\section{Fixed progressions and the Liouville weight}\label{sec:parity}

\subsection{Two residue conditions before division by shared labels}

Fix $M$ and a residue $r\pmod M$, which may be a nonunit. Multiply
the first physical vector by $\ind_{n\equiv r\ (M)}$ and the second
by $\ind_{n'\equiv r+1\ (M)}$. Their norms decrease. The first vector
remains independent of its lists at a fixed position. Thus the general
pairing still bounds the entire residual family.

For a full small prime interval $\mathcal P_g$, the harmonic law
$1/(V_gp)$ is uniform on $(\Z/M\Z)^\times$ to any fixed negative
power of $L$. Partial summation of the fixed-modulus prime estimate
proves this, since $\log p\geq L^{.27}$ and $V_g$ is bounded below.
One shared label from this group suffices. Conditional on every other
label in the independent-label model, multiplication by their product
permutes the units. The collision errors described in
Section~\ref{sec:framework} remain negligible.

Perform this averaging only after the dyads and Mellin integral have
been restored in \eqref{eq:strip}. For a unit $d\pmod M$, the exact
identity is
\begin{equation}\label{eq:orbit}
 \frac1{\varphi(M)}\sum_{t\in(\Z/M\Z)^\times}
 \ind_{tw\equiv r\ (M)}\ind_{t(w+d)\equiv r+1\ (M)}
 =\frac1{\varphi(M)}\ind_{w\equiv rd\ (M)}.
\end{equation}
Subtracting the two congruences forces $td\equiv1$. The unique
possible $t=d^{-1}$ satisfies the first congruence exactly when
$w\equiv rd$. There is no division by $r$ or $w$.
The raw pattern has $d=D_C=1$. Each comparison acquires
$w\equiv rD_C\pmod M$.

To bound those comparisons, first split each free product into its
fixed unit residues modulo $M$. The remaining condition on $w$ is a
fixed residue. With the Fourier convention
$\widehat U(\theta)=\sum_n U(n)\e(\theta n)$,
\[
 \widehat{\ind_{n\equiv s\ (M)}U}(\theta)
  =\frac1M\sum_{j\bmod M}\e(-js/M)\widehat U(\theta+j/M).
\]
Its squared energy is bounded by the corresponding average of the
squared energies. A rational center with denominator $k$ moves to
one with denominator at most $Mk$. Increasing the fixed logarithmic
denominator and width exponents covers all these arcs. In
\eqref{eq:scales}, $H'$ grows faster than every fixed power of $L$
times a small-prime scale, so the new time cutoff is still admissible.
The bounded tests on the free products preserve the bilinear multiplier
estimate.

Apply this argument to the source centered endpoint
$\W(T_x-C_x)$. The source's local energy estimate already holds at
every fixed logarithmic denominator and time exponent. It yields
\begin{equation}\label{eq:unsigned-correlation}
 \sum_{w\equiv r\ (M)}\frac{\Psi(w/x)}w
 \W(w)(T_x(w)-C_x)G(w+1)W_1(w)W_1(w+1)\ll_A L^{-A}.
\end{equation}
We will use this to prove the unsigned progression law. Placing the
residue condition inside $G$ would violate $G(n)=G(n_*)$; the two
physical residue conditions avoid that problem.

\subsection{Ordinary-integer cancellation}

The classical Liouville estimate, in the reduced-residue form recorded
by Pintz \cite[equation (1.5)]{Pintz}, is
\begin{equation}\label{eq:bvlambda}
 \sum_{k\leq R^\theta}\max_{(a,k)=1}
 \left|\sum_{\substack{n\leq R\\n\equiv a\ (k)}}\lambda(n)\right|
 \ll_{A,\theta}R(\log R)^{-A},\qquad \theta<\tfrac12.
\end{equation}
All parameters are fixed. Smooth compactly supported dyadic weights
follow by partial summation, with the sum over moduli inside the
integral of absolute cumulative sums. No supremum over endpoints is
inserted into \eqref{eq:bvlambda}.

If a residue $a\pmod k$ is a nonunit, set $g=(a,k)$ and write
$n=gv$. Then $v\equiv a/g\pmod{k/g}$ is a reduced class and
$\lambda(gv)=\lambda(g)\lambda(v)$. For
$R\geq x^\delta$, $k\leq L^C$, and $|t|\leq L^B$, this gives
\begin{equation}\label{eq:swlambda}
 \frac1R\sum_{n\equiv a\ (k)}\lambda(n)\psi(n/R)n^{it}
 \ll_A L^{-A}.
\end{equation}
The test $\psi$ has fixed compact support in $(0,\infty)$, with
fixed logarithmic derivative costs. The reduced scale $R/g$ is still
a positive power of $x$; increasing the accuracy in
\eqref{eq:bvlambda} pays every displayed logarithmic cost.

\subsection{A signed correlation with a nonempty tuple}

\begin{lemma}\label{lem:lambda-correlation}
For the marked data of Section~\ref{sec:framework}, fixed $M,r$, and
$\Psi\in C_c^\infty((0,\infty))$,
\begin{equation}\label{eq:lambda-correlation}
 \sum_{w\equiv r\ (M)}\frac{\Psi(w/x)}w
 \lambda(w)\W(w)T_x(w)G(w+1)W_1(w)W_1(w+1)
 \ll_A L^{-A}.
\end{equation}
The tuple dimension satisfies $d\geq1$.
\end{lemma}

\begin{proof}
Insert $\lambda(n)$ in the first physical vector. Its norm and
list independence are unchanged. If $s_P$ is the number of active
groups and $c_\nu=\sum_g|C_g|$, then
\[
 \Omega(D_R)=Js_P-c_\nu,\qquad
 \lambda(D_Rw)=(-1)^{Js_P}(-1)^{c_\nu}\lambda(w).
\]
These identities count repeated auxiliary labels with multiplicity.
Divide out the common sign $(-1)^{Js_P}$. The transference result
and \eqref{eq:orbit} give the raw correlation plus comparisons
multiplied by $(-1)^{c_\nu}$, up to $O_A(L^{-A})$.
Each comparison still needs a separate bound.

At low Mellin times $|t|\leq L^B$, expand the tuple and the divisor
model \eqref{eq:divisor}, writing $w=\ell av$.
The scale of $v$ is $R=Y/(\ell a)\geq x^{\eps-o(1)}$.
Complete multiplicativity separates $\lambda(\ell a)$, and splitting
the additive phase into residue classes permits \eqref{eq:swlambda}.
The total normalized error is bounded using
$\sum_{\boldsymbol\ell}1/\ell=O(1)$ and
$\sum|A_a|/a\leq L^C$. Cauchy's inequality controls the divisor-model
error without changing its square norm. Choose the accuracy after
$B$ so that the integral over these low times also saves every
required power of $L$.

At high times remove one small mark and one long prime slot, giving
$w=p_s p_\ell v$. Then
$\lambda(w)=\lambda(v)$, so the small- and long-prime polynomials
are unchanged up to constant signs. Their residual coefficient
has the same square norm as the original coefficient. Small-mark
squares and repeated numerical long-slot primes have square errors
$O(YL^C\exp(-L^{.27}))$ and $O(YL^Cx^{-\eps})$. Ordinary Dirichlet
mean square bounds dispose of these errors on the whole time interval.

The factorial estimate in \cite[Section 6]{D} places the times when
the small-prime polynomial exceeds a prescribed negative power of
$L$ inside at most $\exp(O(L^{.9}))$ unit intervals. Outside these
intervals, ordinary mean square on the collapsed scale $Y/P_s$
applies because $T_0/(Y/P_s)=LP_s/H\to0$. Here $P_s$ denotes the
small-prime scale, so $P_s\leq\exp(L^{.36})$.
On the exceptional intervals, use \eqref{eq:longprime} and the
source sparse Gram estimate for the remaining coefficients. That
estimate is coefficient-independent: its matrix entries are the
unweighted sums $\sum_{v\asymp R}v^{i(t_i-t_j)}$. Multiplying the
tested coefficient by $\lambda(v)$ preserves its square norm.
The long-prime scale lies between $x^\eps/2$ and $x^{1-\eps}$;
the residual scale is at least $x^{\eps-o(1)}$. All retained times
are below $x^2$ after fixed logarithmic twists. These are the source
estimate's ranges. Choose $B$ and the long-prime accuracy after the
fixed comparison costs.

This proves \eqref{eq:energy}. The Fourier shifts for the fixed
residue condition use the same bounds with enlarged fixed exponents.
The local Mellin and multiplier inputs now bound every comparison.
Subtract them from the transferred residual to obtain
\eqref{eq:lambda-correlation}.
\end{proof}

The nonempty tuple supplies the long prime in this proof. We have
not proved a signed constant-endpoint estimate by this method.
The total Liouville mean will follow instead from size-biased sampling.

\subsection{Progression remainders and scalar replacement}

Fix $0<c_0<\eps/8$ and a smooth nonnegative compact cutoff $\Phi$.
For $H$ equal to a fixed inert weight, a product of two such weights,
a constant, or a bounded linear combination of these, the source
represented-divisor expansion and CRT give
\begin{align}\label{eq:unsigned-progression}
 \sum_{\substack{r\leq x^{c_0}\\(r,M)=1}}
 \left|\sum_{\substack{r\mid u+1\\u+1\equiv b\ (M)}}
 \Phi(u/x)H(u)T_x(u)
 -\frac{x}{rM}\left(\int\Phi\right)C_x\E_{\rm div}H\right|
 \ll_A xL^{-A}.
\end{align}
The statement also holds with $T_x,C_x$ replaced by $1,1$.
Indeed every retained tuple/divisor product is eventually coprime
to fixed $M$. A divisor sharing a prime with $r$ gives no solutions;
otherwise CRT gives one class modulo $rM$. The source comparison of
the forbidden inert primes with its independent model has exponentially
small error. The total rounding error is
$x^{1-\eps+c_0+o(1)}$. Divisors with $(r,M)>1$ contribute zero.

The corresponding signed remainder is
\begin{equation}\label{eq:signed-progression}
 \sum_{r\leq x^{c_0}}
 \left|\sum_{\substack{r\mid u+1\\u+1\equiv b\ (M)}}
    \Phi(u/x)H(u)\lambda(u)T_x(u)\right|
 \ll_A xL^{-A}.
\end{equation}
To see this, write $u=\ell av$ in each retained divisor term.
When compatible, the congruences specify a class of $v$ modulo $rM$.
Its gcd with $rM$ is the fixed $g=(b-1,M)$. After $v=gt$, the
class is reduced modulo $rM/g$. The moduli are distinct as $r$ varies
and are less than $(x/(\ell ag))^{1/3}$ for large $x$.
Apply \eqref{eq:bvlambda} at that residual scale and sum with the
reciprocal tuple and divisor coefficients. For the omitted terms use
the positive first-omitted Bonferroni divisor envelope and CRT. Its
factorial reciprocal tail saves every logarithmic power, and its
rounding error has the bound above. This direct progression argument
avoids multiplying an $L^2$ error by $x^{c_0}$.

Finally, the scalar replacement extends to all unit-modulus weights
$f_z(n)=z^{\Omega(n)}$ used here. Selecting one prime from each of
$K$ inert bands and writing $mn-1=ah$ gives
\begin{equation}\label{eq:detphase}
 f_z(mn-1)F_0(mn-1)=z^K f_z(h)F_0(h)
\end{equation}
for an invariant bounded $F_0$. The factor on the right lies outside
the original Cauchy square and has modulus at most one. The scalar
coefficients $\alpha,\beta$ are unchanged, so the assumed coefficient
criterion applies. The selected-prime-square error is absolute.
The product class is imposed by
\[
 \ind_{mn\equiv b\ (M)}
 =\frac1{\varphi(M)}\sum_{\chi\bmod M}
       \overline{\chi(b)}\chi(m)\chi(n).
\]
The character factors preserve \eqref{eq:coefficient}, with logarithmic
test moduli enlarged by a fixed factor. No Liouville or complex weight
is inserted into the prime-minus-rough discrepancy coefficient.

\section{The principal term for a nonreal mode}\label{sec:model}

Fix a nonreal root of unity $z$, and put $f_z(n)=z^{\Omega(n)}$.
The proof below treats each such $z$ separately. All constants may
depend on $z$. In particular, no estimate uniform as $z\to1$ is claimed.

\subsection{Classical distribution inputs}

For $\Rea s>1$,
\begin{equation}\label{eq:euler}
 \sum_{n\geq1}\frac{f_z(n)}{n^s}=\zeta(s)^zG_z(s),\qquad
 G_z(s)=\prod_p\frac{(1-p^{-s})^z}{1-zp^{-s}}.
\end{equation}
Branches are continued from this half-plane. The product $G_z$ is
analytic near $1$. The Selberg--Delange theorem of de la Bret\`eche
and Tenenbaum \cite[Theorem 1.2 and equation (1.10)]{BT} gives a
fixed-function expansion to arbitrary fixed logarithmic precision.
Its hypotheses hold since $f_z(p)=z$, $|f_z|=1$, and the higher
prime-power sums converge with any fixed exponent $\sigma>1/2$.

Set $A_z(s)=((s-1)\zeta(s))^zG_z(s)$. The cumulative coefficients
are those of $A_z(s)/s$, divided by $\Gamma(z-j)$. The smooth
density coefficients are those of $A_z(s)$. This follows by
Stieltjes partial summation of a finite cumulative expansion, with
its remainder bounded in absolute value. In detail, if
$A_z(1+t)=\sum a_jt^j$ and $A_z(1+t)/(1+t)=\sum c_jt^j$, then
$a_j=c_j+c_{j-1}$. Differentiating the finite terms gives
\[
 \frac{c_j}{\Gamma(z-j)}+
 \frac{(z-j)c_{j-1}}{\Gamma(z-j+1)}
 =\frac{a_j}{\Gamma(z-j)}.
\]
There is no differentiation of an unspecified error.

We also require the centered progression estimate
\begin{equation}\label{eq:bvz}
 \sum_{q\leq R^\theta}\max_{(a,q)=1}
 \left|\sum_{\substack{n\leq R\\n\equiv a\ (q)}}f_z(n)
 -\frac1{\varphi(q)}\sum_{\substack{n\leq R\\(n,q)=1}}f_z(n)\right|
 \ll_{z,A,\theta}R(\log R)^{-A},\qquad\theta<\tfrac12.
\end{equation}
Here is the verification of the hypotheses needed to apply
Granville--Shao \cite[Theorem 1.1(b)]{GS}.
For a nonprincipal character $\chi$, the twisted Dirichlet series
is $L(s,\chi)^z$ times a product absolutely convergent in
$\Rea s>1/2$. Assumption~\ref{ass:inputs} gives a zero-free
half-plane $\Rea s>7/8$. On each closed smaller half-plane,
Borel--Carath\'eodory applied to discs centered at $2+it$ bounds the
analytic logarithm of $L$ by $O(\log(q(2+|t|)))$. Polynomial bounds
for $L$ follow from character periodicity and partial summation;
the Euler series fixes the logarithm at the centers. The discs avoid
zeros and poles. Extra Euler factors of imprimitive characters are
nonzero there. Thus the twisted series has polynomial growth.

Smooth Mellin inversion can be moved to $\Rea s=15/16$. For
$q\leq(\log R)^D$, it gives arbitrary logarithmic saving after
choosing a sufficient transform decay order. Smooth majorants with
transition width $R(\log R)^{-K}$ remove the cutoff with an error
smaller than any prescribed negative logarithmic power. Character
orthogonality gives the centered Siegel--Walfisz criterion.
For larger $q$, choose $D$ large and use the elementary bounds
$R/q+1$ and $R/\varphi(q)$, treating $q>R$ separately.

The logarithmic derivative coefficients satisfy
$\Lambda_{f_z}(p^j)=z^j\log p$. Its convolution inverse is
$\mu f_z$, whose logarithmic derivative coefficients are the negatives
of these. Both functions belong to the class required by
\cite{GS}. On primes, $f_z=z\ind_{\mathbb P}$, so classical prime
Bombieri--Vinogradov gives the other hypothesis of Theorem 1.1(b).
This proves \eqref{eq:bvz}. Smooth versions follow with the sum over
moduli inside partial summation and a strict exponent margin at
every comparable scale. The principal branch at $s=1$ has not been
removed by this argument.

\subsection{Definition and exact covariance}

Put $\rho=23/50$, $y=\exp(L^\rho)$, and
\[
 n_y=\prod_{p\leq y}p^{v_p(n)},\qquad
 V_y=\prod_{p\leq y}(1-1/p),\qquad
 P_{z,y}(s)=\prod_{p\leq y}(1-zp^{-s}).
\]
For a sufficiently large fixed integer $N$, define
\begin{equation}\label{eq:density}
 b_{z,y,N}(v)=\sum_{j=0}^{N-1}
 \frac{(A_zP_{z,y})^{(j)}(1)}{j!\Gamma(z-j)}
          (\log v)^{z-1-j},\qquad v>1.
\end{equation}
Choose a fixed smooth real $\chi$, between zero and one, equal to
zero on $(-\infty,\eps/4]$ and one on $[\eps/2,\infty)$.
Set $\beta(v)=\chi(\log v/L)b_{z,y,N}(v)$ for $v>1$ and zero
otherwise. Define the principal term on every positive integer by
\begin{equation}\label{eq:model}
 \Bx(n)=\frac{f_z(n_y)}{V_y}\,z^d
       \sum_{\boldsymbol\ell}\frac1\ell
          \beta\left(\frac{n}{n_y\ell}\right).
\end{equation}
This definition uses all primes at most $y$, not an active subset.
The same function therefore occurs for every active set used later.

\begin{lemma}\label{lem:modelsize}
Let $\delta_z=1-\Rea z>0$. On each fixed range $n\leq x^C$,
\begin{equation}\label{eq:modelsize}
 |\Bx(n)|\ll_{N,z,d,\eps,C}L^{-(1-\rho)\delta_z}.
\end{equation}
For every $y$-smooth integer $a$,
$\Bx(an)=f_z(a)\Bx(n)$. The logarithmic exponents in bounds for
any fixed number of smooth derivatives can be chosen independently
of fixed $N$; their constants and thresholds may depend on $N$.
\end{lemma}

\begin{proof}
The identity $(an)_y=an_y$ gives exact covariance, also when $a$ and
$n$ share prime factors. For each fixed $j\geq1$,
\[
 |(\log P_{z,y})^{(j)}(1)|\ll_{j,z}(\log y)^j,\qquad
 |(A_zP_{z,y})^{(j)}(1)|\ll_{j,z}|P_{z,y}(1)|(\log y)^j.
\]
The second bound also holds for $j=0$.
These estimates follow by differentiation of the finite Euler
product and $\sum_{p\leq y}(\log p)^j/p\ll_j(\log y)^j$.
On the cutoff support, successive density terms have a factor
$O_{j,z,\eps}(L^{\rho-1})$. Also
\[
 \frac{|P_{z,y}(1)|}{V_y}
 =\exp\left(\delta_z\sum_{p\leq y}\frac1p+O_z(1)\right)
 \asymp_z L^{\rho\delta_z}.
\]
Multiply this by $L^{\Rea z-1}$ and the bounded harmonic tuple
sum. This proves \eqref{eq:modelsize}. The same calculation applies
to a fixed number of logarithmic derivatives. The finite sequence
of decreasing powers $L^{-(1-\rho)j}$ absorbs their constants for
large $x$ depending on $N$.
\end{proof}

The leading coefficient of \eqref{eq:density} is
$G_z(1)P_{z,y}(1)/\Gamma(z)$. Taking absolute values of all Euler
coefficients before summing would lose the bound just proved.

\subsection{Truncations with square and progression bounds}

Let $D_0=\exp(L^{1/2})$ and $h_0=\lfloor L^{.01}\rfloor$.
For $\eta=1/\log y$, Rankin's inequality gives
\begin{equation}\label{eq:rankin}
 \sum_{P^+(a)\leq y}\frac{a^\eta}{a}\ll L^C,
 \qquad
 \sum_{\substack{P^+(a)\leq y\\a>D_0/y}}\frac1a
 \ll L^C\exp(-\tfrac12L^{.04}).
\end{equation}
If $n_y>D_0$, multiplying its prime factors one at a time finds a
divisor in $(D_0/y,D_0]$. This gives a positive short-divisor
majorant for the discarded integers. Fixed local divisor moments
remain logarithmic, since their local Euler factors are
$1+O(1/p)+O(1/p^2)$. Cauchy's inequality then gives arbitrary
logarithmic square savings as well. On a represented physical list
with product $P$, the Rankin weight adds at most
$P^\eta\leq L^{O_J(1)}$. The saving in \eqref{eq:rankin} absorbs
that cost after $J$ is fixed.

For roughness use the finite Bonferroni expansion
\begin{equation}\label{eq:bonferroni}
 \ind_{P^-(v)>y}
 =\sum_{\substack{e\mid v, P^+(e)\leq y\\
                    \mu^2(e)=1, \omega(e)\leq h_0}}\mu(e)+E(v),
 \qquad |E(v)|\leq\binom{\omega_{\leq y}(v)}{h_0+1}.
\end{equation}
The last expression is a positive divisor sum. Its divisors are
$\leq\exp(O(L^{.47}))$, and their reciprocal mass is at most
$(C\log L)^{h_0+1}/(h_0+1)!$. The retained absolute reciprocal mass
is $L^{O(1)}$.

For a tuple term $u=\ell v$, a divisor condition $a\mid v$ and
compatible congruences $r\mid u+1$, $u+1\equiv b\pmod M$ give
$O_M(x/(\ell ra)+1)$ integers on the support. There are
$O_d(x^{1-\eps})$ tuples. Every fixed product of the short-divisor
supports above has ordinary coefficient mass $x^{o(1)}$. Summing
the rounding errors over $r\leq x^{c_0}$ therefore costs
\begin{equation}\label{eq:rounding}
 x^{1-\eps+c_0+o(1)}=O_A(xL^{-A}).
\end{equation}
For model terms the tuple is harmonic rather than a divisibility
condition, and the rounding cost is only $x^{c_0+o(1)}$.

For square bounds, sum the omitted roughness terms over the possible
smooth parts before applying a norm inequality. Apart from a fixed
logarithmic factor and the bounded harmonic tuple sum, the pointwise
error is at most
\[
 \tau(n_y)2^{\omega_{\leq y}(n)}
       \ind_{\omega_{\leq y}(n)\geq h_0}.
\]
Squaring, using $4^w\ind_{w\geq h_0}\leq2^{-h_0}8^w$, and applying
a fixed divisor moment gives $O(RL^C2^{-h_0})$ on a dyad of scale $R$.
This is the aggregate square bound needed below. Summing separate
norms with a factor $\sum_{a\leq D_0}a^{-1/2}$ would not suffice.

\subsection{Matching the local principal terms}

\begin{lemma}[Rough density]\label{lem:roughdensity}
Fix a positive residual exponent $\eps_1$ and fixed logarithmic
modulus, time, and derivative bounds. For $R\in[x^{\eps_1},x^C]$,
$k\leq L^C$, $(a,k)=1$, $|t|\leq L^B$, and a smooth dyadic test
$\psi$, every sufficiently large fixed $N$ gives
\begin{equation}\label{eq:roughdensity}
 \frac1R\sum_{\substack{v\equiv a\ (k)\\P^-(v)>y}}
 f_z(v)\psi(v/R)(v/R)^{it}
 =\frac1{\varphi(k)}\int\psi(u)u^{it}b_{z,y,N}(Ru)\dd u
       +O_A(L^{-A}).
\end{equation}
For unweighted rough integers the density is $V_y/\varphi(k)$,
with the same accuracy and smoothness bounds.
\end{lemma}

\begin{proof}
Every prime of $k$ is at most $y$ for large $x$.
Use \eqref{eq:bonferroni}. The nonprincipal character sums save
every power of $L$ by the smooth Siegel--Walfisz estimate proved
above, on scales $R/e\geq x^{\eps_1-o(1)}$.
The retained reciprocal mass and the error envelope pay the fixed
modulus and time costs.

For the principal character, apply the fixed-function
Selberg--Delange expansion to each residual scale $R/e$.
Taylor-expand $\log(Ru/e)$ around $\log(Ru)$. Since
$\log e=O(L^{.47})$, each further Taylor order gains a factor
$L^{-.53}$ before fixed costs. The restored moments are
\[
 \sum_{P^+(e)\leq y}\frac{\mu(e)f_z(e)}e(-\log e)^j
       =P_{z,y}^{(j)}(1).
\]
Restoring the omitted degrees in each fixed moment costs at most
\[
 2^{-h_0}\sum_{e\mid\prod_{p\leq y}p}
       \frac{2^{\omega(e)}(\log e)^j}{e}
 \ll_j 2^{-h_0}L^{C_j}.
\]
The derivative product rule and
\[
 \frac1{\Gamma(z-j)}\binom{z-1-j}{r}
       =\frac1{r!\Gamma(z-j-r)}
\]
give \eqref{eq:density}. Choosing sufficiently many fixed terms
gives the desired error; the bounds in Lemma~\ref{lem:modelsize}
permit any fixed $N$ beyond a sufficient threshold.
The principal character is one on every rough integer, so its
contribution is divided by $\varphi(k)$.
For the unweighted case, CRT in \eqref{eq:bonferroni} gives density
$k^{-1}\prod_{p\leq y,p\nmid k}(1-1/p)=V_y/\varphi(k)$.
\end{proof}

\begin{lemma}[Local matching]\label{lem:localmatch}
Let $H_x(n)=H_x(n_y)$, with $|H_x|\leq L^C$. At $Y=x^{1+o(1)}$,
uniformly for logarithmic $k,t$ and the permitted smooth test,
\begin{equation}\label{eq:localmatch}
 \frac1Y\sum_n\psi(n/Y)(n/Y)^{it}H_x(n)
       [f_z(n)T_x(n)-\Bx(n)]\e(hn/k)=O_A(L^{-A}).
\end{equation}
\end{lemma}

\begin{proof}
Discard $n_y>D_0$ using the preceding bounds. In the genuine tuple
term write $n=\ell av$, where $a=n_y$ and $v$ is $y$-rough.
The coefficient is $z^df_z(a)H_x(a)$, and
$R=Y/(\ell a)\geq x^{\eps-o(1)}$. Lemma~\ref{lem:roughdensity}
gives the principal term
\begin{equation}\label{eq:localmain}
 \frac{z^df_z(a)H_x(a)}{\ell a}
 \frac{c_k(ha\ell)}{\varphi(k)}
 \int\psi(u)u^{it}b_{z,y,N}(Yu/(\ell a))\dd u,
\end{equation}
where $c_k$ is the Ramanujan sum over unit residues.
For the model, write $n=av$ and use unweighted rough counting.
Its $V_y$ cancels the $1/V_y$ in \eqref{eq:model}.
Since $\ell$ is a unit modulo $k$, $c_k(ha\ell)=c_k(ha)$, also
when $a$ is a nonunit. The model term is therefore exactly
\eqref{eq:localmain}. The cutoff $\chi$ is one on these residual
scales. Sum the errors with $\sum_a1/a\leq L^C$ and
$\sum_{\boldsymbol\ell}1/\ell=O(1)$, choosing precision after all
test and time exponents. No regularity of $H_x(a)$ is needed.
\end{proof}

\subsection{The discrete high-time estimate}

Take the first comparison core $f_zT_x-\Bx$. Lemma~\ref{lem:localmatch}
handles every fixed logarithmic range of low times. At high times,
extract one small mark. For the genuine tuple contribution, extract
one long prime as well. The identity
$f_z(p_sp_\ell v)=z^2f_z(v)$ leaves the two prime polynomials
unchanged up to phases. Their residual coefficient has unchanged
square norm, so the proof of Lemma~\ref{lem:lambda-correlation}
applies to this part.

For the model contribution, exact covariance gives
$\Bx(p_sn)=z\Bx(n)$. At the remaining scale
$R_0\asymp Y/P_s$, write $n_y=a\leq D_0$ and expand the roughness
of $n/a$. The resulting terms are
\begin{equation}\label{eq:discretemodel}
 \frac{z^df_z(a)H_{\rm rem}(a)\mu(e)}{V_y\ell}
       \ind_{ae\mid n}\,\beta\left(\frac n{a\ell}\right),
\end{equation}
where $e$ is squarefree, $P^+(e)\leq y$, and $\omega(e)\leq h_0$.
After writing $n=aev$, the normalized reciprocal coefficient mass
is $L^{C_*}$. The smooth amplitude is $\beta(ev/\ell)$; its total
variation obeys the same bound, with exponent independent of fixed $N$.
All residual scales $R'=R_0/(ae)$ satisfy $R'\geq x^{1-o(1)}$.

The omitted coefficient has square mass $O_A(R_0L^{-A})$.
Ordinary Dirichlet mean square, applied before restriction to
exceptional times, gives
\begin{equation}\label{eq:modelerror}
 R_0^{-2}\int_{|t|\leq2T_0}\left|\sum_n E_n n^{it}\right|^2\dd t
 \ll L^{-A}(T_0/R_0+\log R_0),\qquad
 T_0/R_0\asymp LP_s/H=o(1).
\end{equation}
The last relation uses both lower bounds on $U_i$ in
\eqref{eq:scales}, not only $Y=x^{1+o(1)}$.

Let $T_{\rm mod}=\exp(L/(\log L)^2)$. For
$L^B\leq|t|\leq T_{\rm mod}$, separate the additive phase into
residue classes modulo $k$. Sum--integral comparison and partial
summation on each progression give a normalized bound
\[
 |N_0(t)|\ll L^{C_*}(|t|^{-1}+x^{-c}).
\]
Indeed the discrete error is $O(L^C(1+|t|)/R')$, while integration
by parts gives the $|t|^{-1}$ term. Consequently
\begin{equation}\label{eq:moderate}
 \int_{L^B\leq|t|\leq T_{\rm mod}}|N_0(t)|^2\dd t
 \ll L^{2C_*-B}+x^{-c'}.
\end{equation}
Logarithmic twists are absorbed by increasing $B$.
For $T_{\rm mod}<|t|\leq2T_0$, the exact ranges in
\eqref{eq:logphase} apply at $R'$ and give
\begin{equation}\label{eq:highmodel}
 |N_0(t)|\ll L^{C_*}\exp(-L/(\log L)^{C_{19}}).
\end{equation}
The residue may be a nonunit; there is no inversion of $ae$.

Outside the exceptional small-prime intervals, the small polynomial
and ordinary mean square on scale $Y/P_s\gg T_0$ suffice.
On the exceptional intervals, integrate \eqref{eq:moderate}
directly and use \eqref{eq:highmodel} over at most
$\exp(O(L^{.9}))$ unit intervals. The latter loss is absorbed by
its exponential saving. This proves \eqref{eq:energy} for the
model part. Choose $B$ after the fixed kernel, group, and derivative
exponents, then choose $N$ after $B$ and the low-time accuracy.
The derivative exponents in Lemma~\ref{lem:modelsize} make this
order possible.

\begin{proposition}[Marked model transfer]\label{prop:modeltransfer}
For the data of Section~\ref{sec:framework}, every fixed $M,r$ and
every required accuracy $A$, all sufficiently large fixed $N$ satisfy
\begin{equation}\label{eq:modeltransfer}
 \sum_{w\equiv r\ (M)}\frac{\Psi(w/x)}w\W(w)
 [f_z(w)T_x(w)-\Bx(w)]G(w+1)W_1(w)W_1(w+1)
 \ll_A L^{-A}.
\end{equation}
One $N$ can be used for every active set in a fixed extraction.
\end{proposition}

\begin{proof}
The first physical core is eventually bounded, so
\eqref{eq:listnorm} gives a norm exponent independent of $J,N$.
It remains list-independent. Use the original ideal family.
On stripping shared labels, both terms have the exact covariance
\[
 [f_zT_x-\Bx](D_Rw)
      =z^{Js_P-c_\nu}[f_zT_x-\Bx](w).
\]
Divide out $z^{Js_P}$, restore the dyads and Mellin integral, and
apply \eqref{eq:orbit}. The residual contains the raw difference
and all comparison differences with factors $z^{-c_\nu}$.
The local and high-time estimates just proved bound every comparison,
including fixed residue restrictions through Fourier shifts.
Choose their accuracy after the total logarithmic pattern cost and
subtract them. The same $\Bx$ and a sufficiently large single $N$
work for every active set, because their exponents have a common
fixed bound in each extraction.
\end{proof}

\section{Progression remainders and extraction of primes}\label{sec:extraction}

\subsection{The complex progression remainder}

Retain fixed $M,b$ with $(b,M)=1$, $0<c_0<\eps/8$, and a smooth
nonnegative compact cutoff $\Phi$. We first prove the estimate needed
by the positive sieve.

\begin{proposition}\label{prop:complexremainder}
For the inert-array weights $H$ allowed in
\eqref{eq:unsigned-progression}, every $A>0$, and sufficiently large
fixed $N$,
\begin{equation}\label{eq:complexremainder}
 \sum_{\substack{r\leq x^{c_0}\\(r,M)=1}}
 \left|\sum_{\substack{r\mid u+1\\u+1\equiv b\ (M)}}
    \Phi(u/x)H(u)[f_z(u)T_x(u)-\Bx(u)]\right|
 \ll_A xL^{-A}.
\end{equation}
\end{proposition}

\begin{proof}
Here $H(u)=H(u_y)$ and is bounded for each fixed array collection.
Discard $u_y>D_0$ by the positive divisor envelope
\eqref{eq:rankin} and the summed rounding estimate
\eqref{eq:rounding}. In a genuine tuple term write
$u=\ell av$, where $a=u_y\leq D_0$ and $P^-(v)>y$.
The long product $\ell$ is a unit modulo $rM$ for large $x$.
Set $g=(b-1,M)$. Compatibility requires
\begin{equation}\label{eq:compatibility}
 (a,r)=1,\qquad (a,M)=g.
\end{equation}
When these conditions hold, $v$ lies in a reduced class modulo
$q=rM/g$. Prime by prime, partial removal of a prime-power factor
of $M$ fixes the valuation of $a$; full removal allows larger
valuations. In each case $a/g$ is a unit modulo $M/g$.
Roughness still excludes every prime at most $y$, including those
fully removed from the modulus. The model term $u=av$ has the same
compatibility and modulus, although its reduced residue can differ.

Insert \eqref{eq:bonferroni} in the genuine progression sum and in
its centered natural mean. Only divisors $e$ coprime to $q$ can
contribute in a reduced class. Their residual scales satisfy
$R/e\geq x^{\eps-o(1)}$, and the distinct moduli $rM/g$ are
less than $(R/e)^{1/3}$. Apply \eqref{eq:bvz}, then sum with the
$L^{O(1)}$ reciprocal divisor mass. The positive envelope controls
progression errors by CRT; natural-mean errors cost only a
logarithmic reciprocal-totient sum.

In the remaining natural means, the exclusion of primes $p\mid q$
with $p>y$ can be removed. After division by $\varphi(q)$, its
total cost is bounded using
\begin{align}\label{eq:largeq}
 \sum_{q\leq Mx^{c_0}}\frac1{\varphi(q)}
     \sum_{\substack{p\mid q\\p>y}}\frac1p
 &\leq\sum_{p>y}\frac1{p(p-1)}
           \sum_{t\leq Mx^{c_0}/p}\frac1{\varphi(t)}
 \ll L^{O(1)}/y.
\end{align}
We used $\varphi(pt)\geq(p-1)\varphi(t)$, including $p\mid t$.
This saves every logarithmic power. Lemma~\ref{lem:roughdensity}
at modulus one now gives, for fixed $a,\ell$, principal contribution
\begin{equation}\label{eq:presievemain}
 \frac{xz^df_z(a)H(a)}{a\ell\varphi(q)}
    \int\Phi(t)b_{z,y,N}(xt/(a\ell))\dd t.
\end{equation}

For the model, use unweighted rough counting in the reduced class
of $v$. Its density is
\begin{equation}\label{eq:roughq}
 \frac{V_y}{q\prod_{p\mid q, p\leq y}(1-1/p)}
 =\frac{V_y}{\varphi(q)}\prod_{p\mid q, p>y}(1-1/p).
\end{equation}
The $1/V_y$ in \eqref{eq:model} cancels $V_y$.
Equation~\eqref{eq:largeq} permits replacement of the final product
by one. The result is exactly \eqref{eq:presievemain}, including
every retained density coefficient. These terms cancel before
absolute values are taken. Sum the errors with the reciprocal
$a$-mass $L^{O(1)}$ and tuple mass $O(1)$, using
\eqref{eq:rounding} for the integer-counting errors.
\end{proof}

\subsection{A common positive sieve argument}

We now treat three first-endpoint differences:
\begin{equation}\label{eq:threecores}
 D_x(u)=
 \begin{cases}
 T_x(u)-C_x,&\text{unsigned progression law},\\
 \lambda(u)T_x(u),&\text{parity},\\
 f_z(u)T_x(u)-\Bx(u),&\text{nonreal mode}.
 \end{cases}
\end{equation}
Their summed progression remainders vanish to arbitrary logarithmic
precision by \eqref{eq:unsigned-progression},
\eqref{eq:signed-progression}, and
Proposition~\ref{prop:complexremainder}. For the unsigned case,
subtract $C_x$ times the remainder without the tuple.

Choose a large fixed even $h$, put $\gamma=c_0/(4h+3)$, and sieve
$u+1$ by primes at most $x^\gamma$ that do not divide $M$.
For the unsigned case apply the block sieve separately to
$\Phi\W T_x$ and $C_x\Phi\W$. For parity use the nonnegative
weights $\Phi\W T_x(1\pm\lambda)/2$.
For the complex case, $|\Bx|\leq1$ eventually, so with
$A_0=\Phi\W(T_x+1)$ the four weights
\[
 A_0\pm\Phi\W\Rea D_x,\qquad
 A_0\pm\Phi\W\Ima D_x
\]
are nonnegative. All these weights include the same progression
indicator. Each plus/minus pair has identical main terms and
arbitrary-precision summed remainders. The block sieve is applicable
because $x^{\gamma(4h+2)}<x^{c_0}$. Also
\[
 \frac1M\prod_{p\leq x^\gamma, p\nmid M}(1-1/p)
 =\frac1{\varphi(M)}\prod_{p\leq x^\gamma}(1-1/p)
\]
for large $x$. Subtracting the paired estimates gives
\begin{equation}\label{eq:sifted}
 \limsup_{x\to\infty}\frac{L}{xm_{\W}}
 \left|\sum_{\substack{u+1\equiv b\ (M)\\P^-(u+1)>x^\gamma}}
       \Phi(u/x)\W(u)D_x(u)\right|
 \ll_{M,d,\eps,\Phi}e^{-h}/\gamma.
\end{equation}
The surviving constant is independent of $K,h$ and the particular
fixed inert array. Array-dependent remainder constants vanish in
the first limit. In the complex case, the unsigned upper sieve and
\eqref{eq:modelsize} bound the presieved model by
$o(xm_{\W}/L)$ for fixed $h$ and array. Thus
\eqref{eq:sifted} also holds with $D_x$ replaced by $f_zT_x$.

\subsection{Replacing composite factors and retaining the coin coefficients}

A presieved composite in the support has a fixed bounded number
$j$ of prime factors, with $2\leq j\leq2/\gamma$ for large $x$.
Repeated large factors have negligible absolute mass.
Telescope each prime factor to the positive rough proxy of
Section~\ref{sec:framework}. The scalar criterion and
\eqref{eq:detphase} permit this for the unsigned, parity, and
complex statistics. Fixed characters impose the product residue.
Clipping, coefficient truncation, dyadic splitting, and logarithmic
separation have the source's original absolute error bounds.
Choose $K$ after these fixed losses and $\gamma$, independently of
the later inert exponents.

Write $v=\prod_{i=1}^jm_i$, $u=v-1$, and let $\mathcal C$ be the
full candidate set of small and big groups. A bar on $m_i$ removes
all candidate prime parts, including inactive groups. The stripped
coefficient is
\[
 \widetilde a(m_i)=\ind_{P^-(m_i)>W}
     \ind_{\bar m_i>x^\gamma}\kappa(\bar m_i),
 \qquad \kappa(t)=\frac1{V(W)\log t}.
\]
The source bounds its total positive tuple mass by
$O_{j,\gamma}(x/L)$ and the threshold/logarithm replacement error
by $o(x/L)$.

Conditional on the tuple, the source uses independent coins with
probabilities
\begin{equation}\label{eq:coins}
 s_g(u,v)=c_1q_0^{\omega_g(u)-1}\frac{\omega_g(u)}{V_g}
              (q_0/j)^{\omega_g(v)-1}\frac{\omega_g(v)}{V_g}.
\end{equation}
Choose $c_1>0$ so that these numbers lie in $[0,1]$; a zero hit
count gives zero. Let $\mathcal R$ be the success sets with at least
$c_2\log L$ groups of each kind. The source's success-probability
estimate says that the other outcomes have total positive proxy
mass $o(x/L)$. Discard them at this stage, before expanding signs.
All first-endpoint statistics under consideration are bounded.

The retained-event probability has the exact expansion
\begin{align}\label{eq:coinexpand}
 R(u,v)
 &=\sum_{S\in\mathcal R}\prod_{g\in S}s_g
                     \prod_{g\notin S}(1-s_g)\notag\\
 &=\sum_{S\in\mathcal R}\sum_{T\subseteq\mathcal C\setminus S}
                  (-1)^{|T|}\prod_{g\in S\cup T}s_g,
 \qquad 0\leq R(u,v)\leq1.
\end{align}
Each active set $P=S\cup T$ meets the hypotheses of the marked
correlation, and the number and coefficient mass of terms are
$L^{O(1)}$.

Outside candidate-prime-square exceptions, strip the active part
of $v$, obtaining $v_*$. For each factorization $e_1\cdots e_j=v_*$,
the active primes can be allocated independently to the $j$ slots.
The bars make the coefficients identical for all allocations.
The identity
\begin{equation}\label{eq:allocation}
 j^{\omega_P(v)}(q_0/j)^{\omega_P(v)-|P|}
       =j^{|P|}q_0^{\omega_P(v)-|P|}
\end{equation}
then gives the invariant core
\[
 G_P(v)=\sum_{e_1\cdots e_j=v_*}
 \prod_{i=1}^j\left[
  \ind_{P^-(e_i)>W}\ind_{\bar e_i>x^\gamma}\kappa(\bar e_i)\right].
\]
It satisfies $G_P(v)=G_P(v_*)$ and
$|G_P(v)|\leq L^{C_j}\tau(v_*)^{C_j}$.
Every expanded term has its actual coefficient
$(-1)^{|T|}c_1^{|P|}j^{|P|}$ multiplying the marked sum.
Candidate-square exceptions have exponentially small divisor-weighted
mass and remain negligible after this logarithmic expansion cost.

For the unsigned and parity cases, apply
\eqref{eq:unsigned-correlation} or
Lemma~\ref{lem:lambda-correlation}, with $r=b-1$ and
$\Psi(t)=t\Phi(t)$, and multiply by $x$.
Choose the accuracy after the expansion costs. This makes the
composite contribution $o(x/L)$.

For a nonreal mode, Proposition~\ref{prop:modeltransfer} replaces
the first core by $\Bx(u)$ with arbitrary logarithmic precision.
The principal contribution for a given $P$ is
\[
 \sum_{u+1\equiv b\ (M)}\Phi(u/x)\W(u)\Bx(u)
                W_{1,P}(u)W_{1,P}(u+1)G_P(u+1).
\]
Sum these terms with all their coefficients before estimating their
absolute values. Reversing \eqref{eq:allocation} and then
\eqref{eq:coinexpand} gives exactly
\begin{equation}\label{eq:aggregate}
 \sum_{m_1,\ldots,m_j}
 \ind_{v\equiv b\ (M)}\Phi((v-1)/x)\W(v-1)\Bx(v-1)
       R(v-1,v)\prod_i\widetilde a(m_i),
\end{equation}
up to the already controlled square exceptions.
The same $\Bx$ occurs for every active set, so it factors through
both identities. All other factors in \eqref{eq:aggregate} are
nonnegative, and $R\leq1$. Hence its absolute value is at most
\begin{equation}\label{eq:aggregatesmall}
 \sup_{u\asymp x}|\Bx(u)|\,O_{j,K,\gamma,\Phi}(x/L)=o(x/L).
\end{equation}
The source mass bound follows from $O(xV(W)^j)$ rough tuples per
product box, coefficient size $O_\gamma((LV(W))^{-j})$, and
$O(L^{j-1})$ boxes. The fixed inert weight costs a fixed constant.
Logarithmic expansion factors multiply only the arbitrary-precision
errors, not \eqref{eq:aggregatesmall}. This distinction permits roots
$z$ whose fixed positive exponent in \eqref{eq:modelsize} is small.

\subsection{Removing the inert marks}

We have bounded the prime sum with each fixed inert array by the
right side of \eqref{eq:sifted}. For any fixed integer $B\geq1$,
choose $B$ disjoint arrays of the same sufficient size $K$, and put
\[
 Z_B(u)=\left(\frac1B\sum_{a=1}^B\frac{\W_a(u)}{m_a}-1\right)^2.
\]
The source nonnegative deweighting estimate is
\begin{equation}\label{eq:deweight}
 \limsup_{x\to\infty}\frac Lx\sum_p\Phi((p-1)/x)Z_B(p-1)
 \ll C_K/B.
\end{equation}
It applies without a progression restriction. Restriction decreases
its left side. The other Cauchy factor is $O((x/L)^{1/2})$, because
$T_x-C_x$, $\lambda T_x$, and $f_zT_x$ are bounded on the support.
Consequently the unweighted statistics satisfy
\begin{equation}\label{eq:limsup}
 \limsup_{x\to\infty}\frac Lx
 \left|\sum_{p\equiv b\ (M)}\Phi((p-1)/x)F_x(p-1)\right|
 \leq C_{M,d,\eps,\Phi}\frac{e^{-h}}\gamma
                    +C_{d,\eps,\Phi}\sqrt{C_K/B},
\end{equation}
where $F_x$ is respectively $T_x-C_x$, $\lambda T_x$, or $f_zT_x$.
The model term has already been removed in the last case.

The order of limits is part of the proof. Fix the progression and
interior data, then $h,\gamma$, the finite range of $j$, and a
sufficient $K$. Fix any finite number $B$ of arrays. Choose the
physical families, high-time exponents, density order $N$, and error
precisions. Take $x\to\infty$ first, then $B\to\infty$, and finally
$h\to\infty$, choosing a new sufficient fixed $K$ for each $h$.
Since $\gamma^{-1}=(4h+3)/c_0$, \eqref{eq:limsup} tends to zero.
Constants depending on $N,K,B,\gamma$ in errors have disappeared in
the first limit.

\begin{theorem}[Interior estimates]\label{thm:interior}
Under Assumption~\ref{ass:inputs}, for every fixed reduced class
$b\pmod M$, every nonempty interior tuple statistic
\eqref{eq:tuple}, and every $\Phi\in C_c^\infty((0,\infty))$,
\begin{align}
 \sum_{p\equiv b\ (M)}\Phi((p-1)/x)(T_x(p-1)-C_x)&=o(x/L),
       \label{eq:interiorunsigned}\\
 \sum_{p\equiv b\ (M)}\Phi((p-1)/x)z^{\Omega(p-1)}T_x(p-1)&=o(x/L)
       \quad(z\ne1),\label{eq:interiorsigned}
\end{align}
where $z$ is any fixed root of unity different from one.
\end{theorem}

The proof used $\Phi\geq0$ for the positive sieve. A real smooth
compact cutoff is a difference of two such cutoffs, obtained with a
smooth dominating bump. Real and imaginary parts then give the
complex formulation by linearity.

\section{Probability limits and squarefree restrictions}\label{sec:probability}

\subsection{The unsigned law in each progression}

\begin{proposition}\label{prop:unsigned}
Under Assumption~\ref{ass:inputs}, $V(p)$ converges in each fixed
finite set of coordinates to $\PD(1)$ when $p$ ranges over a fixed
reduced progression, in either the dyadic or cumulative convention.
\end{proposition}

\begin{proof}
First take $p$ uniform in $(x,2x]$ in the class. The denominator
$N_M(x)$ is asymptotic to $x/(\varphi(M)L)$ by the fixed-modulus
prime number theorem. For rectangles $\prod_{i=1}^d(a_i,b_i]$ with
positive lower endpoints and $\sum b_i<1$, choose prime slots
$(x^{a_i},x^{b_i}]$. Their reciprocal tuple sum tends to
$\prod_i\log(b_i/a_i)$; coincidences contribute $o(1)$.
Equation~\eqref{eq:interiorunsigned} and smooth approximation of
the dyadic cutoff give the same limit for the mean tuple count.

Regard the prime factors of $p-1$, with multiplicity, as labelled
balls of masses $\log q/\log(p-1)$. On a compact subset of
\[
 \mathcal D_d=\{(t_1,\ldots,t_d):t_i>0,\ \sum_i t_i<1\},
\]
the expected ordered distinct-ball measure therefore converges to
$\prod_i\dd t_i/t_i$. The change from $\log x$ to $\log(p-1)$
is $1+O(1/L)$ and is controlled by endpoint strips. Repeated
numerical primes on a fixed interior compact set entail
$q^2\mid p-1$ with $q\geq x^{a/2}$ for some fixed $a>0$.
Their integer count is
$O(x^{1-a/2}+\sqrt{x})=o(N_M(x))$. Tuple counts are bounded on
such a compact set, so multiplicities cause no further error.

Sample balls without replacement with probability proportional to
their remaining masses. The first $d$ size-biased draws multiply
the factorial density by
\[
 \prod_{i=1}^d\frac{t_i}{1-\sum_{j<i}t_j}.
\]
The resulting density is
$h_d(t)=\prod_{i=1}^d(1-\sum_{j<i}t_j)^{-1}$ on $\mathcal D_d$.
The triangular substitution
$u_i=t_i/(1-\sum_{j<i}t_j)$ identifies it with uniform measure
on $(0,1)^d$. Its mass is one. Compact exhaustion therefore excludes
escaped boundary mass and proves convergence of the size-biased
draws to the first $d$ $\GEM(1)$ fragments.

Let $R_{x,d}$ be the unseen mass after $d$ draws, with zero padding
after exhaustion. Its expectation tends to $2^{-d}$.
Sorting the first $d$ draws and padding by zero approximates every
ranked coordinate of $V(p)$ within $R_{x,d}$.
Uniform continuity and Markov's inequality, followed by $d\to\infty$,
prove the ranked finite-dimensional law.
A finite decomposition of $(X/2^J,X]$ into dyads gives the cumulative
law. The omitted initial class-prime proportion tends to $2^{-J}$;
take $X\to\infty$ and then $J\to\infty$.
\end{proof}

\subsection{Complex measures and Fourier inversion}

Attach $f_z(p-1)$ to each sampled prime, for a fixed $z\ne1$.
The resulting complex factorial measure $\sigma_{x,d}^{z}$ has
total variation dominated by its unsigned counterpart.
Equation~\eqref{eq:interiorsigned} implies
\[
 \int\phi\dd\sigma_{x,d}^{z}\longrightarrow0,
       \qquad\phi\in C_c(\mathcal D_d).
\]
The same unsigned endpoint-strip and multiplicity bounds used above
justify the normalization and rectangle-to-continuous-test passage
in absolute value.

Let $\nu_{x,d}$ be the unsigned size-biased law and
$\nu_{x,d}^{z}$ its complex weighted version. On the interior,
\[
 \dd\nu_{x,d}^{z}
  =\prod_i\frac{t_i}{1-\sum_{j<i}t_j}\dd\sigma_{x,d}^{z},
 \qquad |\nu_{x,d}^{z}|\leq\nu_{x,d}.
\]
Choose $0\leq\chi\leq1$ compactly supported in $\mathcal D_d$
with $\int\chi h_d>1-\eta$. For any continuous $H$ on the cube,
\[
 \left|\int H\dd\nu_{x,d}^{z}\right|
 \leq\left|\int\chi H\dd\nu_{x,d}^{z}\right|
       +\|H\|_\infty\int(1-\chi)\dd\nu_{x,d}.
\]
The first term tends to zero, and the limsup of the second is at
most $\eta\|H\|_\infty$. Let $\eta\to0$. Thus
$\nu_{x,d}^{z}$ converges weakly to the zero complex measure on
the whole cube. This is weak convergence, not convergence in total
variation. At $d=1$ its total mass is exactly $\E f_z(p-1)$,
which proves cancellation without a separate $d=0$ correlation.

If $H$ tests $k$ ranked coordinates, the error from sorting $d$
draws, with $d\geq k$, is at most
\[
 \omega_H(\eta)+2\|H\|_\infty\E R_{x,d}/\eta.
\]
The same bound holds for the complex weighted average because
$|f_z|=1$. Take $x\to\infty$, then $d\to\infty$, then
$\eta\to0$. We obtain
\begin{equation}\label{eq:modecancel}
 \E[f_z(p-1)H(V_1(p),\ldots,V_k(p))]\longrightarrow0
       \quad(z\ne1).
\end{equation}
For fixed $m$, finite Fourier inversion gives
\[
 \ind_{\Omega(n)\equiv a\ (m)}
  =\frac1m\sum_{j=0}^{m-1}
       \exp(-2\pi ija/m)\exp(2\pi ij\Omega(n)/m).
\]
The trivial mode is Proposition~\ref{prop:unsigned}; $z=-1$ is
the parity case; the other modes were treated in
Section~\ref{sec:model}. This proves Theorem~\ref{thm:main} in
dyads. The same finite dyadic decomposition proves its cumulative
form.

\subsection{Local square exclusions}

Let $B_Y(n)$ indicate that $q^2\nmid n$ for every prime $q\leq Y$.
Within a fixed reduced class modulo $M$, this condition on $p-1$
is a finite union of reduced classes modulo
$\lcm(M,\prod_{q\leq Y}q^2)$, apart from finitely many primes.
For fixed $Y$, apply the unsigned law and \eqref{eq:modecancel}
with $z=-1$ in each class. Reduced-lift counting gives the factors
$1-\rho_q(M,b)$ in \eqref{eq:rho}.
If $q\nmid M$, there are $q(q-1)$ reduced lifts modulo $q^2$
and one has $p\equiv1\pmod{q^2}$. If $v_q(M)=1$ and
$b\equiv1\pmod q$, one of $q$ lifts has that property. If
$q^2\mid M$, it is forced or excluded by $b$ itself.

To let $Y$ grow after $x$, use
\begin{equation}\label{eq:squaretail}
 \limsup_{x\to\infty}
 \frac{\#\{x<p\leq2x:\exists q>Y, q^2\mid p-1\}}
      {\pi(2x)-\pi(x)}\ll1/Y.
\end{equation}
For $Y<q\leq x^{1/4}$, Brun--Titchmarsh (in the form recorded in \cite{Ford}) gives
$O(x/[q(q-1)\log x])$ primes in the relevant progression.
For $q>x^{1/4}$, count integers and use
\[
 \sum_{x^{1/4}<q\leq\sqrt{2x}}(x/q^2+1)
       =O(x^{3/4}+\sqrt{x})=o(x/\log x).
\]
The first reciprocal tail is $O(1/Y)$.
Changing the denominator in \eqref{eq:squaretail} to $N_M(x)$
costs only a fixed factor. Bounded tests and signs preserve the bound.

Write $B=\mu^2$ temporarily for the arithmetic indicator.
Taking $x\to\infty$ and then $Y\to\infty$ yields the four limits
\begin{align*}
 \E H(V(p))&\longrightarrow\E H(L),&
 \E[\lambda(p-1)H(V(p))]&\longrightarrow0,\\
 \E[B(p-1)H(V(p))]&\longrightarrow C_{M,b}\E H(L),&
 \E[\lambda(p-1)B(p-1)H(V(p))]&\longrightarrow0.
\end{align*}
The identities
\[
 \ind_{\lambda=s,B=1}=B(1+s\lambda)/2,\qquad
 \ind_{\lambda=s,B=0}=(1-B)(1+s\lambda)/2
\]
give every joint atom. This proves Theorem~\ref{thm:squarefree},
including $C_{M,b}=0$. Finally $\mu=\lambda\mu^2$ gives
\eqref{eq:mu} and the two M\"obius-sign densities.

\subsection{The prescribed-valuation consequence}

On the event $v_q(p-1)=e_q$ for $q\in\mathcal S$, let
$B_{\mathcal S}$ indicate squarefreeness outside $\mathcal S$.
Then
\[
 \mu_{\mathcal S}(p-1)
       =(-1)^{\sum_{q\in\mathcal S}e_q}\lambda(p-1)B_{\mathcal S}(p-1).
\]
For fixed $Y$, the valuation event and the square exclusions at
$q\leq Y$, $q\notin\mathcal S$, are a finite union of reduced
classes modulo
\[
 \lcm\left(\prod_{q\in\mathcal S}q^{e_q+1},
           \prod_{\substack{q\leq Y\\q\notin\mathcal S}}q^2\right).
\]
Apply \eqref{eq:modecancel} with $z=-1$ and $H=1$ in each class.
Equation~\eqref{eq:squaretail} controls the remaining square factors.
Taking the limits in that order proves Corollary~\ref{cor:moree}.
The argument permits $e_q\geq2$, because squarefreeness is imposed
only outside $\mathcal S$.

\appendix
\section{Analytic dependencies}\label{sec:sources}

The results are conditional on the following inputs, with the hypotheses
and normalizations stated in Section~\ref{sec:framework}. References
\cite{D,S,P,Z} refer throughout to the same fixed repository revision.\footnote{Revision
\texttt{fd4aeeb2ee4fc729c18d98444fed42fd0529eeeb} of
\href{https://github.com/openai/math/tree/fd4aeeb2ee4fc729c18d98444fed42fd0529eeeb}{\texttt{openai/math}}.
The bibliography links to these versions.}
Proof estimates mentioned below are included among the assumptions.

\subsection*{Transference and marked correlations}
We use the general graph-transference and residual ideal-family
propositions in \cite[Section 6]{D}, which apply the local transference,
endpoint pairing, and ideal-kernel results of \cite[Sections 3--4]{S}.
The required endpoint norms, list normalization, shared-label collision
bound, and bilinear multiplier estimate are those used in their proofs.
The local Mellin-energy lemma, truncated divisor-model lemma, and sparse
Gram bound in \cite[Section 6]{D} give the reduction
\eqref{eq:energy} and expansion \eqref{eq:divisor}.
In particular, the norm exponent is independent of the fixed shared-list
parameter $J$. Sections~\ref{sec:parity} and~\ref{sec:model} verify the
new endpoint energy estimates; no signed or fixed-progression version
of the final distribution theorem is assumed.

\subsection*{Oscillation and the scalar estimate}
The long-prime polynomial and logarithmic-phase lemmas in
\cite[Appendix A]{D} are used in the ranges
\eqref{eq:longprime}--\eqref{eq:logphase}, including nonunit residue
classes in the latter. For the determinant replacement, we use the
coefficient criterion for the marked estimate in \cite[Section 8]{P}
and the scalar Cauchy-square bound in its proof. Its coefficient
hypothesis is \eqref{eq:coefficient}; the prime-minus-rough coefficient
satisfies it by the minor-moment and major-term estimates in
\cite[Sections 3--5]{D}. Equation~\eqref{eq:detphase} keeps the new
multiplicative factor outside that scalar square.

\subsection*{Positive extraction}
We use the block-sieve and divisor-moment estimates in
\cite[Section 2]{D}, followed by the progression, prime-to-rough,
conditional-coin, allocation, and mean-square deweighting estimates in
\cite[Section 7]{D}. These include the proof estimates for removing
candidate prime parts from the rough coefficients and thresholds.
The sieve level, positive proxy mass, and deweighting bound used here
are stated in Sections~\ref{sec:framework} and~\ref{sec:extraction}.
The complex progression remainder and the recombination of its
principal terms are proved in Section~\ref{sec:extraction}.

\subsection*{Ordinary multiplicative-function estimates}
The main theorem of \cite{Z} supplies the zero-free half-plane for
Dirichlet $L$-functions required in Section~\ref{sec:model}. We use
this assertion as a hypothesis. The remaining distribution inputs are
the Liouville estimate recorded in \cite[equation (1.5)]{Pintz},
Granville--Shao's criterion \cite[Theorem 1.1(b)]{GS}, and the
Selberg--Delange expansion \cite[Theorem 1.2 and equation (1.10)]{BT},
as applied in Sections~\ref{sec:parity} and~\ref{sec:model}.

\begingroup
\small
\begin{thebibliography}{99}
\bibitem{D}
OpenAI,
\emph{The Poisson--Dirichlet law for prime predecessors}, 2026.
\href{https://github.com/openai/math/tree/fd4aeeb2ee4fc729c18d98444fed42fd0529eeeb/preprints/The-Poisson-Dirichlet-Law-for-Prime-Predecessors-September-24-2026/build}{Source manuscript}.

\bibitem{S}
OpenAI,
\emph{Weighted dilation graphs, smooth shifted primes, and totient fibers}, 2026.
\href{https://github.com/openai/math/tree/fd4aeeb2ee4fc729c18d98444fed42fd0529eeeb/preprints/Weighted-Dilation-Graphs-Smooth-Shifted-Primes-and-Totient-Fibers-September-24-2026/build}{Source manuscript}.

\bibitem{P}
OpenAI,
\emph{Primitive roots for every admissible integer base}, 2026.
\href{https://github.com/openai/math/tree/fd4aeeb2ee4fc729c18d98444fed42fd0529eeeb/preprints/Primitive-roots-for-every-admissible-integer-base-October-4-2026/build}{Source manuscript}.

\bibitem{Z}
OpenAI,
\emph{The quasi-Riemann hypothesis}, 2026.
\href{https://github.com/openai/math/tree/fd4aeeb2ee4fc729c18d98444fed42fd0529eeeb/preprints/The-Quasi-Riemann-Hypothesis-September-30-2026/build}{Source manuscript}.

\bibitem{BT}
R. de la Bret\`eche and G. Tenenbaum,
\emph{Remarks on the Selberg--Delange method},
Acta Arith. \textbf{200} (2021), 349--369.
\href{https://arxiv.org/abs/2010.12929v8}{arXiv:2010.12929v8}.

\bibitem{Carella}
N. A. Carella,
\emph{Result on the Mobius function over shifted primes}, 2022.
\href{https://arxiv.org/abs/2206.12956v3}{arXiv:2206.12956v3}.

\bibitem{Ford}
K. Ford,
\emph{An explicit sieve bound and small values of $\sigma(\varphi(m))$}.
\href{https://www.ford126.web.illinois.edu/wwwpapers/sigphi.pdf}{Author's manuscript}, introduction.

\bibitem{GS}
A. Granville and X. Shao,
\emph{When does the Bombieri--Vinogradov theorem hold for a given
multiplicative function?}, 2017.
\href{https://arxiv.org/abs/1706.05710v1}{arXiv:1706.05710v1}.

\bibitem{Lichtman}
J. D. Lichtman,
\emph{Averages of the M\"obius function on shifted primes}, 2021.
\href{https://arxiv.org/abs/2009.08969v2}{arXiv:2009.08969v2}.

\bibitem{MH}
P. Moree and H. Hommersom,
\emph{Value distribution of Ramanujan sums and of cyclotomic
polynomial coefficients}, 2003.
\href{https://arxiv.org/abs/math/0307352v1}{arXiv:math/0307352v1}.

\bibitem{Pintz}
J. Pintz,
\emph{Are there arbitrarily long arithmetic progressions in the
sequence of twin primes? II}, 2010.
\href{https://arxiv.org/abs/1004.1067v1}{arXiv:1004.1067v1}.
\end{thebibliography}

\endgroup
\end{document}
